The multi-seed confirmatory evidence covers two datasets, RetinaMNIST and
Solar, under CE- and RPS-trained representations. This provides a focused
cross-setting test of the frozen intervention, but it does not establish the
same response for other ordinal tasks, data modalities, backbones, objectives,
or imbalance regimes. UTKFace contributes provenance-clean historical seed-0
evidence only. Its seeds 1--4 are unavailable, so it is supplementary
supporting evidence and is not part of the primary H1/H2a confirmatory block.

The two confirmatory populations also have different roles. RetinaMNIST uses a
training-only OOF head evaluation, whereas Solar uses a fixed archived readout.
The Solar readout supports the predeclared backbone-realization comparison, but
it is not a new external cohort for each seed. Endpoint supports and dataset
structures also differ, so the results should not be interpreted as directly
comparable estimates of localization severity across datasets.

The representation analysis is deliberately limited. Train-derived centroid
distances provide a coarse diagnostic of endpoint-specific geometry; they do
not identify a causal representation failure or determine whether an individual
sample is recoverable. This limitation is reflected in the H2a outcome: the
geometry diagnostic is strong in Solar, partial in RetinaMNIST CE, and not
supported in RetinaMNIST RPS. The paper therefore does not claim a universal
geometry-conditioned recovery mechanism.

The controlled RetinaMNIST factor study is also narrow. It evaluates one frozen
RPS representation and finds that balanced sampling is the dominant observed
adaptation signal in that setting; it does not establish a general causal role
for balanced sampling. The separate support-severity grid shows that lower
class-4 support weakens direct endpoint evidence, but its localization outcomes
are non-monotonic. These results do not causally identify imbalance severity as
the source of inward localization.

Finally, direction-only adaptation is an explanatory probe rather than a
production method. It shows that changing classifier direction can alter
rare-end localization under the fixed protocol, but it does not show that the
head is the sole source of error, that all rare-end samples are recoverable, or
that the intervention improves global accuracy, calibration, risk, or
uncertainty quality. A practical correction would require separate evaluation
of its broader predictive and deployment consequences. This scope remains
intentional.
